\documentclass[conference]{IEEEtran}
\IEEEoverridecommandlockouts

\usepackage{cite}
\usepackage{makecell}
\usepackage{amsmath,amssymb,amsfonts}
\usepackage{mathtools}
\usepackage{tabularx}
\usepackage{booktabs}
\usepackage{graphicx}
\usepackage[utf8]{inputenc}
\usepackage[brazilian]{babel}
\usepackage[hidelinks]{hyperref}
\usepackage{textcomp}
\usepackage{accsupp}
\usepackage[T1]{fontenc}
\usepackage{amsmath}            
\usepackage{float}              
\usepackage{algorithm}           
\usepackage{algpseudocode}       
\usepackage{listings}
\usepackage{xcolor}
\def\BibTeX{{\rm B\kern-.05em{\sc i\kern-.025em b}\kern-.08em
    T\kern-.1667em\lower.7ex\hbox{E}\kern-.125emX}}
% --- Anexo de uso de IAG ---
\newcounter{usoiag}
\newcommand{\usoIAG}[4]{%
  \refstepcounter{usoiag}%
  \par\medskip\noindent\hrule\smallskip
  \noindent\textbf{Uso \theusoiag}\\
  \textbf{Etapa / local de uso:} #1\\
  \textbf{Modelo utilizado:} #2\\
  \textit{\textbf{Prompt:}} \textit{#3}\\
  \textbf{Observações:} #4\par
}

\floatname{algorithm}{Algoritmo}
\algrenewcommand\algorithmicrequire{\textbf{Entrada:}}
\algrenewcommand\algorithmicensure{\textbf{Saída:}}
\algrenewcommand\algorithmicfor{\textbf{para}}
\algrenewcommand\algorithmicdo{\textbf{faça}}
\algrenewcommand\algorithmicif{\textbf{se}}
\algrenewcommand\algorithmicthen{\textbf{então}}
\algrenewcommand\algorithmicelse{\textbf{senão}}
\algrenewcommand\algorithmicreturn{\textbf{retorne}}


\definecolor{codebg}{RGB}{245, 245, 245}
\definecolor{codegreen}{RGB}{0, 128, 0}
\definecolor{codegray}{RGB}{128, 128, 128}
\definecolor{codepurple}{RGB}{128, 0, 128}
\definecolor{codeblue}{RGB}{0, 0, 180}


\newcommand*{\nonselectablenumber}[1]{{\BeginAccSupp{ActualText={}}#1\EndAccSupp{}}}


\lstdefinestyle{codigo}{
    backgroundcolor=\color{codebg},
    commentstyle=\color{codegreen},
    keywordstyle=\color{codeblue},
    numberstyle=\tiny\color{codegray}\nonselectablenumber,
    stringstyle=\color{codepurple},
    basicstyle=\ttfamily\scriptsize,
    breakatwhitespace=false,
    breaklines=true,
    captionpos=b,
    nolol=false,
    keepspaces=true,
    numbers=right,
    numbersep=5pt,
    showspaces=false,
    showstringspaces=false,
    showtabs=false,
    tabsize=3,
    frame=single,
    framerule=0.5pt,
    rulecolor=\color{codegray},
    extendedchars=true,
    literate=
    {á}{{\'a}}1 {ã}{{\~a}}1 {â}{{\^a}}1 {à}{{\`a}}1
    {é}{{\'e}}1 {ê}{{\^e}}1 {è}{{\`e}}1
    {í}{{\'i}}1 {î}{{\^i}}1 {ì}{{\`i}}1
    {ó}{{\'o}}1 {ô}{{\^o}}1 {õ}{{\~o}}1 {ò}{{\`o}}1
    {ú}{{\'u}}1 {ü}{{\"u}}1 {ù}{{\`u}}1
    {ç}{{\c{c}}}1 {Ç}{{\c{C}}}1
}

\lstset{style=codigo}

\begin{document}

\title{Controle Inteligente de Nível de Água em um Reservatório através com o Microcontrolador ATmega328P}
\author{
\IEEEauthorblockN{Gabriel Luduvice Nunes Wolfovitch, João Pedro Nascimento Monteiro e \\
João Vítor de Araújo Costa de Oliveira}
\IEEEauthorblockA{\textit{Departamento de Engenharia Elétrica e Computação (DEEC)}\\
\textit{Universidade Federal da Bahia (UFBA)}\\
Salvador, Brasil\\
\{gabrielwolfovitch, joaovaco, monteirojoao\}@ufba.br
}
}

\maketitle

\begin{abstract}
This document is a model and instructions for \LaTeX.
This and the IEEEtran.cls file define the components of your paper [title, text, heads, etc.]. *CRITICAL: Do Not Use Symbols, Special Characters, Footnotes, 
or Math in Paper Title or Abstract.
\end{abstract}

\begin{IEEEkeywords}
component, formatting, style, styling, insert.
\end{IEEEkeywords}

\section{Introdução}
This document is a model and instructions for \LaTeX.
Please observe the conference page limits. For more information about how to become an IEEE Conference author or how to write your paper, please visit   IEEE Conference Author Center website: https://conferences.ieeeauthorcenter.ieee.org/.

\section{Parametrização do Problema}
% Seção para descrever os parâmetros escolhidos, como a quantidade de modos amostrados, a capacidade do tanque, dentre outros.

\section{Módulos Utilizados}
%Seção para descrever cada uma das implementações feitas, como o adc, o temporizador e a usart.


\subsection{Conversor Analógico-Digital (ADC)}
\subsection{Comunicação Serial (USART)}
\subsection{Temporizador}

\section{Análise dos Resultados }

\section{Conclusão}

\appendices

\section{Uso de Inteligência Artificial Generativa}

\subsection*{Declaração}
A equipe declara que a IAG foi utilizada nas etapas descritas abaixo. O detalhamento de cada uso, incluindo os \textit{prompts} empregados, é apresentado a seguir.

\subsection*{Detalhamento}

\usoIAG
  {Demonstração Matemática do Algoritmo e Erros Associados}
  {\textit{Claude Sonnet} 5.5 \textit{(Anthropic)}}
  {Demonstre matematicamente o funcionamento do algoritmo dígito por dígito em ponto fixo, considerando formatação Q16.0 na entrada e Q8.6 na saída. Prove também qual o limitante superior dos erros envolvidos}
  {}

\begin{thebibliography}{00}
\bibitem{b1} G. Eason, B. Noble, and I. N. Sneddon, ``On certain integrals of Lipschitz-Hankel type involving products of Bessel functions,'' Phil. Trans. Roy. Soc. London, vol. A247, pp. 529--551, April 1955.
\end{thebibliography}


\end{document}
